\begin{description}
\item[(a)] For the Lagrange points L1 and L2, use the following diagram.
\ioaafig{2025_TH_Q08-s1.pdf}
For Mars in its orbit around the Sun, $F_{\mathrm{cf}}=F_{SM}$,
\[
M_M\omega^2r_{MS}=G\frac{M_MM_\odot}{r_{MS}^2}.
\]\hfill(2 points)
Hence $\omega=\sqrt{GM_\odot/r_{MS}^3}$. \hfill(1 point)
For a mass $m$ at L1,
\[
m\omega^2(r_{MS}-x_{L1})=G\frac{mM_\odot}{(r_{MS}-x_{L1})^2}-G\frac{mM_M}{x_{L1}^2}.
\]\hfill(2 points)
Substituting $\omega^2$ and rearranging,
\begin{align*}
\frac{M_\odot}{r_{MS}^3}(r_{MS}-x_{L1})&=\frac{M_\odot}{(r_{MS}-x_{L1})^2}-\frac{M_M}{x_{L1}^2},\\
M_\odot\left(1-\frac{x_{L1}}{r_{MS}}\right)^3&=M_\odot-M_M\left(\frac{r_{MS}}{x_{L1}}-1\right)^2.
\end{align*}
Since $x_{L1}\ll r_{MS}$, use
\[
\left(1-\frac{x_{L1}}{r_{MS}}\right)^3\approx1-3\frac{x_{L1}}{r_{MS}}
\]\hfill(2 points)
and $r_{MS}/x_{L1}-1\approx r_{MS}/x_{L1}$. \hfill(1 point)
Thus
\[
3M_\odot\frac{x_{L1}}{r_{MS}}=M_M\left(\frac{r_{MS}}{x_{L1}}\right)^2,\qquad
x_{L1}^3=\frac{M_M}{3M_\odot}r_{MS}^3,
\]
so $x_{L1}=r_{MS}\sqrt[3]{M_M/(3M_\odot)}$. \hfill(2 points)
Numerically,
\[
x_{L1}=1.524\times1.496\times10^{11}\sqrt[3]{\frac{0.107M_\oplus}{3\times333054M_\oplus}}
=1.08\times10^9\,\mathrm m=1.08\times10^6\,\mathrm{km}.
\]\hfill(2 points)
The ratio is $d_{\mathrm{Phobos}M}/x_{L1}\approx9\times10^{-3}$. \hfill(1 point)
\item[(b)] For a mass $m$ at L2,
\[
m\omega^2(r_{MS}+x_{L2})=G\frac{mM_\odot}{(r_{MS}+x_{L2})^2}+G\frac{mM_M}{x_{L2}^2}.
\]\hfill(2 points)
As before,
\[
M_\odot\left(1+\frac{x_{L2}}{r_{MS}}\right)^3=M_\odot+M_M\left(\frac{r_{MS}}{x_{L2}}+1\right)^2.
\]
For $x_{L2}\ll r_{MS}$,
\[
\left(1+\frac{x_{L2}}{r_{MS}}\right)^3\approx1+3\frac{x_{L2}}{r_{MS}}
\]\hfill(2 points)
and $r_{MS}/x_{L2}+1\approx r_{MS}/x_{L2}$. \hfill(1 point)
Therefore $3M_\odot x_{L2}/r_{MS}=M_M(r_{MS}/x_{L2})^2$, yielding
\[
x_{L2}=r_{MS}\sqrt[3]{\frac{M_M}{3M_\odot}}.
\]\hfill(2 points)
For Mars--Phobos,
\[
x_{L2}=d_{\mathrm{Phobos}M}\sqrt[3]{\frac{M_{\mathrm{Phobos}}}{3M_M}}
=9380\sqrt[3]{\frac{1.5\times10^{-8}M_M}{3M_M}}\approx16\,\mathrm{km}.
\]\hfill(2 points)
L2 is close to the surface of Phobos, whose radius is 12.5\,km.
\item[(c)]
\ioaafig{2025_TH_Q08-s2.pdf}
The tidal force is the difference between the attraction of Mars on mass $m$ at the satellite's near surface and on the same mass at its centre. Writing $d=d_{\mathrm{Phobos}M}$ and $R=R_{\mathrm{Phobos}}$,
\[
F_{\mathrm{tidal}}=G\frac{mM_M}{(d-R)^2}-G\frac{mM_M}{d^2}.
\]\hfill(4 points)
For $R\ll d$, $(1-R/d)^{-2}\approx1+2R/d$. \hfill(2 points)
Consequently
\[
F_{\mathrm{tidal}}=G\frac{mM_M}{d^2}\left(1+\frac{2R}{d}-1\right)
=2G\frac{mM_MR}{d^3}.
\]\hfill(1 point)
Phobos's attraction is $F_{\mathrm{grav}}=GmM_{\mathrm{Phobos}}/R^2$. \hfill(1 point)
The Roche limit satisfies $F_{\mathrm{tidal}}=F_{\mathrm{grav}}$. \hfill(1 point)
Therefore
\[
d_{\mathrm{Roche}}^3=2\frac{M_M}{M_{\mathrm{Phobos}}}R^3,\qquad
d_{\mathrm{Roche}}=R\sqrt[3]{2\frac{M_M}{M_{\mathrm{Phobos}}}}.
\]\hfill(2 points)
Numerically, $d_{\mathrm{Roche}}=12.5\sqrt[3]{2/(1.5\times10^{-8})}=6386\,\mathrm{km}$. \hfill(2 points)
The inward distance is $\Delta d=9380-6386=2994\,\mathrm{km}$. \hfill(1 point)
The required time is
\[
\Delta t=\frac{\Delta d}{v}=\frac{2994000}{1.8\times10^{-2}}\,\mathrm{yr}=1.67\times10^8\,\mathrm{yr}.
\]\hfill(1 point)

\textbf{Alternative solution.} Model Phobos as two spheres, each of mass $M_{\mathrm{Phobos}}/2$ and density equal to that of Phobos. \hfill(3 points)
Their radius follows from $2(4\pi r^3/3)=4\pi R^3/3$, giving $r=R/\sqrt[3]2$. \hfill(1 point)
The mutual gravitational force is
\[
F_{\mathrm{grav}}=G\frac{M_{\mathrm{Phobos}}^2}{4(2r)^2}.
\]\hfill(2 points)
The difference in the force due to Mars is
\[
\Delta F=G\frac{M_MM_{\mathrm{Phobos}}}{2(d-r)^2}-G\frac{M_MM_{\mathrm{Phobos}}}{2(d+r)^2}.
\]\hfill(2 points)
Applying the approximation,
\begin{align*}
\Delta F&\approx G\frac{M_MM_{\mathrm{Phobos}}}{2d^2}
\left(1+\frac{2r}{d}-1+\frac{2r}{d}\right)\\
&=G\frac{2M_MM_{\mathrm{Phobos}}r}{d^3}.
\end{align*}\hfill(2 points)
Equating the forces gives
\[
G\frac{2M_MM_{\mathrm{Phobos}}r}{d_{\mathrm{Roche}}^3}=G\frac{M_{\mathrm{Phobos}}^2}{4(2r)^2}.
\]\hfill(2 points)
The printed key continues with
\[
\frac{32r^3}{d_{\mathrm{Roche}}^3}=\frac{M_{\mathrm{Phobos}}}{M_M},\qquad
d_{\mathrm{Roche}}=2R\sqrt[3]{2\frac{M_M}{M_{\mathrm{Phobos}}}}.
\]\hfill(1 point)
This gives $d_{\mathrm{Roche}}=12772\,\mathrm{km}$. \hfill(1 point)
On this model Phobos is already inside the Roche limit. \hfill(1 point)
The official key awards full credit for this alternative model.
\item[(d)] Let $d_S$ be the distance from the Sun's centre to the centre of mass, and $d$ the distance from the centre of mass to L4.
\ioaafig{2025_TH_Q08-s3.pdf}
The diagram gives
\[
h^2=d_{MS}^2-\frac{d_{MS}^2}{4},\qquad h=1.32\,\mathrm{au},\qquad l=2h=2.64\,\mathrm{au}.
\]\hfill(3 points)
\end{description}
